%!TEX root = ../hutbthesis_main.tex

\chapter{基于模块化架构的人车模拟器开发插件总体设计}

\section{系统设计原则与目标}

\subsection{系统设计原则}

为了保证插件系统具有高稳定性、高可维护性、强扩展性、良好的用户体验，本文在总体设计阶段严格遵守五大设计原则，在需求分析、架构设计、模块划分、接口定义、流程设计的各个环节中都贯彻了这五大设计原则。

模块化设计原则就是把复杂的系统按照功能边界分成独立的、低耦合的、高内聚的功能模块，每一个模块只负责一个职责，方便独立开发、单元测试、版本迭代、后期维护，从而降低系统的复杂度，提高开发效率。

可扩展性原则，把API定义、调试模板、打包配置、用户自定义规则全部用JSON格式的配置文件来表示，核心逻辑和配置完全分离，可以快速适应新的SDK版本、新增模拟器（CARLA/AirSim）、增加新的检测规则，不需要修改核心代码，提高系统的通用性。

易用性原则是以人车模拟器开发者为中心，简化操作流程、降低使用门槛，把复杂的配置自动化、高频的操作快捷化、重要的信息可视化，使初学者快速入门、高手高效工作。

实时性原则，代码补全、错误检测、状态反馈、参数更新都在毫秒级完成，不会影响用户的输入，也不会对编辑的流畅性造成影响，保证沉浸式的开发体验。

跨平台兼容性原则把Windows系统列为首选平台，并且可以同时兼容MacOS以及Linux系统，在使用时考虑团队合作、教学实验以及跨设备转移等各方面的考虑因素来保证系统的通用性。

\subsection{系统设计目标}

本文主要设计并实现一个面向人车模拟器、使用VS
Code作为开发工具、包含整个开发链路的快速开发验证插件系统，从根本上解决传统开发模式效率低、调试难、验证周期长、环境配置复杂、监控能力弱等痛点，主要的设计目标如下所示

创建领域专用智能编码辅助系统，按照HUTB、MCP的专属API给出正确的代码补全、即时语法语义校对、参数提示、文档悬浮、弃用标注等功效，削减记忆花费，削减拼写失误，加快编码速度，使编码速度提升40\%以上。

第二，具备一键式多场景调试功能，自带单脚本、主程序、多机通信、远程附加这四种调试模板，自动生成配置、注入环境变量、匹配Python解释器，彻底摆脱手动编写launch.json、反复调试配置的繁琐过程，大大缩减了调试准备时间。

第三，创建仿真的热负载参数设置方式，在运行期间可以改变车速、传感器、天气、光照等参数，不必重新启动仿真环境就可以进行参数的迭代调节，参数迭代的反应速度为0.1秒之内，大幅度缩减了验证时间和开发的效率。

第四，创建一体化的便携式部署方案，把VS Code、Python、插件、SDK、示例脚本、配置文件一起打包成免安装、免配置、解压即用的绿色环境，消除环境差异、简化团队协作、降低教学部署成本。

在此基础上开发出具有实时性特点的可视化监控功能，即可以把车辆状态、传感器数据、仿真指标等信息以可视化的方式呈现出来，支持数据导出、离线演示、异常报警等功能，使开发人员可以迅速找到并解决问题。

第六，在自动化批量测试时可以对不同的参数组合以及不同的测试用例同时执行，产生测试结果并制作成可视化的报告，使测试范围得到扩充，开发开支得以减少，从而使软件质量得到保证。

最终系统要具有高稳定性、高准确率、良好的通用性、强扩展性、跨平台适配性，成为人车模拟器开发领域高效、可靠、易用的专用开发工具。

\section{系统总体架构设计}

1.整体架构

系统采用分层式模块化架构，从上到下分为表现层、业务层和数据层三部分，层与层之间存在单向依赖关系，职责明确，耦合度低，功能完整且具有良好的扩展性，总体架构如图\ref{fig:3-1}所示。

表现层是直接面对用户交互的部分，使用 VS Code
的原生界面体系，把状态栏、命令面板、右键菜单、错误提示面板、WebView
监控面板、弹窗提示等交互组件集成在一起，接收用户的操作，显示系统的状态，反馈执行的结果，从而达到直观、方便、高效的交互效果。

业务层为系统的主要功能实现层，它起到表现层和数据层之间的桥梁作用，主要有代码智能补全、实时错误检测、仿真参数热重载、多场景调试、实时监控、批量测试、环境打包、状态栏交互等8大模块的功能，主要是完成数据处理、逻辑计算、任务调度、协议通信等一系列重要的业务逻辑。

数据层为系统数据支撑层，负责统一存储、管理、加载系统运行所用的各种数据资源，即HUTB/MCP
API定义文件、调试模板配置、打包脚本、用户自定义规则、缓存数据、示例脚本等，给业务层提供稳定的、可靠的、可配置的数据支持。

分层架构一方面可以保证系统具有高的内聚性、低耦合性，另一方面也为以后的功能扩展、模块升级以及跨平台适配打下了灵活的架构基础。

\begin{figure}[H]
\centering
\begin{center}
\includegraphics[width=0.86\textwidth,height=0.58\textheight,keepaspectratio]{images/image4.png}
\end{center}
\caption{系统分层架构图}
\label{fig:3-1}
\end{figure}

根据分层架构，业务层又被细分为八个独立的功能模块，各个模块功能单一、职责分明、互相配合、互不影响，模块之间用标准化的接口进行通信，模块架构如图\ref{fig:3-2}所示。

代码智能补全模块，即当用户输入 hutb./mcp.
前缀的时候，可以提供对应的API补全、文档悬浮、代码片段生成、弃用标记等功能。

实时错误检测模块，对静态语法语义进行校验，对未知函数进行识别，对参数进行校验，对相似度进行推荐，并给出错误提示。

仿真参数热重载模块，文件监听、参数解析、合法性校验、MCP协议通信、结果反馈。

多场景调试模块主要对模板管理、配置生成、环境变量注入、解释器匹配、远程调试进行调试。

实时监控面板模块主要是遥测数据的拉取、可视化渲染、数据导出、离线演示、异常告警。

批量测试模块主要是对测试配置进行解析，执行用例并统计结果，生成报告和分析指标。

环境打包模块是对于资源进行复制、配置、编译脚本、打包归档和部署的总称。

状态栏交互模块主要是对环境状态进行展示，快速执行命令，实时反馈，显示信息提示。

十三大部分一起构成了系统的功能矩阵，就编码、校验、调试、监控、测试、部署全过程的工作进行了描述。

\begin{figure}[H]
\centering
\begin{center}
\includegraphics[width=0.86\textwidth,height=0.58\textheight,keepaspectratio]{images/image5.png}
\end{center}
\caption{系统功能模块图}
\label{fig:3-2}
\end{figure}

2.技术选型

为保证系统稳定、兼容、开发高效、扩展性好，根据VS
Code生态及人车模拟器开发场景要求，选择如下。

开发语言为
TypeScript，强类型、静态检查、编译期报错，可以有效地减少运行时的异常，符合
VS Code 扩展开发规范。

运行环境使用的是VS Code
1.80及以上版本，跨平台、轻量、高扩展、生态成熟，给插件提供稳定运行的环境。

核心框架：VS Code Extension API（界面、命令、语言服务）、Debug Adapter
Protocol（调试）、WebView+ECharts（可视化）、MCP 协议（通信）。

关键依赖库：archiver（ZIP
打包）、fs-extra（文件操作）、path（路径处理）、axios（网络通信）、echarts（数据可视化）、debugpy（Python
调试）。

数据格式为JSON（API定义、配置、模板、测试用例），简洁易读易维护，兼容性好。

技术选型要考虑成熟度、适配性、开发效率、后续扩展，给系统的稳定运行和持续迭代提供技术支持。

\section{核心模块设计}

\subsection{代码智能补全模块}

代码智能补全模块是提高人车模拟器脚本编写效率的基础模块，主要针对HUTB、MCP两个专属API系统，解决接口数量多、命名规则复杂、参数繁杂、容易忘记、频繁查文档效率低等明显不足。模块使用事件驱动、实时匹配、静态缓存的思想，事先解析并结构化存储所有的API信息，在用户编码的时候以毫秒级的响应速度给出精准的智能提示。

模块内部被划分成触发识别、API加载、补全生成、代码片段、弃用标记这五个子单元。触发识别子单元实时监听编辑器输入，准确识别出
hutb.、mcp. 命名空间前缀，过滤无关输入，防止无效弹窗；API
加载子单元在插件初始化的时候一次性读取并解析 JSON 格式的 API
定义文件，创建全局函数映射字典，缓存所有的接口名称、参数结构、功能说明、使用示例、弃用状态，减少重复解析，提高响应速度；补全生成子单元根据匹配结果动态生成标准化补全条目，包含函数签名、功能摘要、参数说明、返回值类型；代码片段子单元自动生成可以直接填充的代码模板，支持
Tab
键快速切换参数占位符，减少手动输入；弃用标记子单元识别废弃接口并做视觉区分，推荐最新替代函数，引导开发者规范编码。

整体设计考虑到了精准度、实时性、易用性、规范性等各方面的要求，可以大大减轻
API 记忆的负担，减少拼写错误，提高编码的流畅度，是高效开发的基础。

\subsection{实时错误检测模块}

实时错误检测模块是保证代码质量、及时拦截运行时异常的重要模块，对于Python这种动态语言语法松散、编译期校验缺失、错误滞后暴露的特性，采用静态分析、语法解析、语义校验、相似度匹配相结合的综合检测方法，在编码阶段就对代码进行扫描，提前发现并标记出各种潜在错误，大大降低了调试的成本，缩短了开发的时间。

模块内部又分为语法解析、函数匹配、参数校验、相似度推荐、诊断渲染五个子单元。语法解析子单元使用抽象语法树（AST）技术，逐行解析Python代码，构建结构化的语法树，准确地提取出函数调用、变量引用、赋值操作等关键的语法节点，函数匹配子单元使用正则表达式准确地匹配出hutb.xxx()、mcp.xxx()这样的调用方式，过滤掉非目标函数，参数校验子单元使用状态机参数计数算法，智能地区分出正常代码、字符串、注释、多层嵌套括号，准确地统计出参数的数量，并且校验出必填项是否完整、参数格式是否合规，相似度推荐子单元使用Levenshtein编辑距离算法，计算未知函数和标准API之间的相似度，自动推荐最接近合法接口，有效解决拼写错误，诊断渲染子单元把检测结果转换成编辑器波浪线、侧边面板提示，标注错误位置、说明原因、给出修正建议，交互直观清晰。

模块包含未知函数、参数缺失、参数超限、弃用接口、调用顺序、格式错误等诸多高频问题，检测准确率高，反应速度快，可以很好地削减掉70\%的运行时错误，明显改善代码的规范程度和稳定性。

\subsection{多场景调试模块}

多场景调试模块可以简化调试的环境，适应复杂的开发环境，对于人车模拟器的开发来说，调试环境多种多样，配置复杂，环境变量很多，解释器匹配困难，远程调试困难，所以内置了四种标准调试模板，一键配置、一键启动、全自动环境适配，极大降低调试门槛、提高调试效率。

模块内部可以分为模板管理、配置生成、环境注入、解释器匹配、远程适配这五个子单元。模板管理子单元预封装了单脚本调试、模拟器主程序调试、多机通信调试、远程附加调试四种标准模板，可以覆盖从简单的脚本到复杂的多机协同的所有调试需求，配置生成子单元会解析模板参数，自动生成
VS Code 标准 launch.json 配置文件，不需要手动编写
JSON，防止出现格式错误，环境注入子单元能够自动识别并配置
HUTB\_HOME、PYTHONPATH、MCP\_HOST、MCP\_PORT
等重要环境变量，保证脚本可以正确加载
SDK、建立通信、识别资源，解释器匹配子单元会根据优先级自动匹配本地 Python
解释器，保证开发、调试、运行环境完全一致，不会因为依赖不同而产生冲突，远程适配子单元可以支持本地和远程路径映射、进程附加、断点调试、变量查看，适合多机协同、硬件在环、远程仿真等复杂的场景。

模块设计既具通用性又具备自动化特性，并且保持稳定以及方便的操作，完全消除了繁杂的调试设置，让开发者可以把更多时间投入到逻辑检验当中去，从而大幅改善了调试效率与开发进程的速度。

\subsection{仿真参数热重载模块}

仿真参数热重载模块属于提高迭代速度、缩减验证周期的新颖模块，针对传统模式下参数改变必须停机、重新启动仿真，迭代速度非常慢的缺陷，依托MCP通信协议达成运行时参数即时更新、秒级生效、无须重启的目的，支持车速、传感器、天气、光照、控制指令等多种参数的实时调节，大幅加快了验证的速度。

模块内部又可以分为文件监听、参数解析、合法性校验、协议通信、结果反馈五个子单元。文件监听子单元实时监听Python脚本保存事件，准确识别热重载专属函数调用，参数解析子单元提取函数参数值、类型、名称，形成更新数据包，合法性校验子单元内置参数范围规则，校验车速（0～250km/h）、传感器采样率（1～240Hz）、视场角（1～180°）等关键参数，阻止非法输入，协议通信子单元利用MCP协议向仿真引擎发出更新指令，采取异步通信来保证实时性，结果反馈子单元把更新成功或者失败，更新的内容以及错误的原因及时地显示在状态栏上，交互及时清楚。

模块具有毫秒级响应、高成功率、强可靠性的特点，可以大大提高参数迭代效率几十倍，是实现快速验证、高频迭代的重要保证。

\subsection{实时监控面板模块}

实时监控面板模块属于可视化展示仿真状况、迅速找到异常的重要模块，就传统模式下仿真数据零散、缺少直观展现、异常难以察觉、剖析繁杂等痛点而言，依托WebView+ECharts技术创建可视化监控界面，及时显示车辆、传感器、环境主要数据，支持趋势剖析、数据导出、离线演示，让开发者直观掌控仿真运行情况，迅速找出问题。

模块内部又分为数据拉取、可视化渲染、交互控制、数据导出、离线模拟五个子单元。数据拉取子单元定时从仿真引擎遥测接口获取车速、加速度、方向盘转角、激光点数、位置坐标、帧率等主要数据；可视化渲染子单元把数据转化成指标卡片、趋势曲线、二维轨迹图，实时平滑刷新；交互控制子单元支持图表缩放、视图切换、告警高亮，交互友好；数据导出子单元一键导出CSV格式的历史数据，便于离线分析；离线模拟子单元在没有仿真环境的时候生成模拟数据，可以进行功能演示和学习。

模块设计考虑实时性、直观性、交互性、实用性，给仿真调试、异常排查、效果分析提供重要的可视化支持。

\subsection{批量测试模块}

批量测试模块属于自动化验证、保证代码质量的关键模块，针对传统测试人工繁杂、覆盖率低、统计容易出错、效率低下的缺点，达成多参数组合、多场景并行、自动执行、自动统计、自动报告的一体化测试功能，大幅提高测试效率和质量。

模块内部又分为配置解析、参数注入、用例执行、结果采集、报告生成五个子单元。配置解析子单元读取JSON格式的测试配置，解析场景、天气、速度、传感器参数等多组测试用例；参数注入子单元用环境变量自动传递参数，不需要修改脚本；用例执行子单元批量运行仿真脚本，自动捕获执行日志；结果采集子单元解析日志，提取完成率、碰撞次数、平均速度、耗时等核心指标；报告生成子单元自动生成HTML可视化的报告，包含概览、用例详情、指标对比图表、通过率统计等内容，直观地展示出测试结果。

模块具有自动化属性，可以自定义，可视化以及可复用，从而大大提高测试的覆盖范围，并且能够有效地减少人工工作量，保证代码质量的好坏，进而改善软件的迭代稳定性。

\subsection{环境打包模块}

环境打包模块是解决环境差异、简化部署、提高协作效率的重要模块，针对传统模式下环境配置复杂、缺少依赖、版本不一致、团队协作困难、教学部署繁琐等问题，使用绿色便携、一体化打包方案，把VS
Code、Python、插件、SDK、示例、配置、脚本统一整合成一个免安装、免配置、解压即用的便携包，完全消除环境差异。

模块内部又可以分为资源复制、配置生成、脚本编写、压缩归档、分发适配这五个子单元。资源复制子单元完全复制VS Code便携版、Python解释器、插件、HUTB/MCPSDK、示例脚本等全部依赖；配置生成子单元自动生成VS Code设置、调试配置、环境变量配置，解压后即可配置完成；脚本编写子单元生成Windows、macOS、Linux启动脚本，一键启动完整的开发环境；压缩归档子单元使用高效的压缩算法生成ZIP包，体积小；分发适配子单元支持跨平台打包，适合多系统部署。

模块实现了全链路的自动化、便携化、跨平台化，大幅降低了部署成本，提高了协同效率，并且保证了环境的一致性。

\subsection{状态栏交互模块}

状态栏交互模块属于轻量化交互、提高便捷性的辅助模块，针对传统模式下状态不透明、高频操作繁琐、反馈不及时等问题，依靠
VS Code
底部状态栏，轻量化集成状态展示和快捷操作，低侵入、高便捷，改善整体使用体验。

模块内部被划分成状态采集、界面渲染、快捷触发、实时反馈、初始化这五个子单元。状态采集子单元实时获取Python版本、SDK连接状态、插件运行状态、热重载结果等重要信息，界面渲染子单元用简单的图标和文字来显示，不占用编辑区，快捷触发子单元绑定点击事件，一键启动调试、打包、刷新API、打开示例，实时反馈子单元显示操作结果、错误提示、进度，初始化子单元在插件启动时完成状态加载和界面初始化。

模块设计简单高效，不侵入，可以即时反应，在不影响正常开发的同时提高交互的方便性和状态的透明度。

\subsection{场景快照与断点续跑模块}

场景快照与断点续跑模块是面向长流程仿真、异常中断与多轮实验场景设计的关键模块，用于解决传统仿真执行中断后需从头运行、中间状态不可回溯、实验复现效率低的问题。模块支持仿真运行状态的捕获、序列化存储、加载恢复与断点续跑，有效提升仿真实验的可追溯性与执行效率。

模块内部划分为快照捕获、状态持久化、快照加载、断点续跑、快照管理五个子单元。快照捕获子单元实时解析当前仿真脚本执行上下文，自动提取车辆动力学状态、传感器配置参数、环境场景信息、热重载参数及当前代码执行行号；状态持久化子单元将捕获的结构化数据序列化为
JSON
格式，统一存储于工作区.hutb/snapshots目录；快照加载子单元读取指定快照文件，解析并还原仿真运行上下文；断点续跑子单元动态更新调试配置文件，注入快照路径、执行行号及相关环境变量，使仿真从指定断点继续执行；快照管理子单元提供快照命名、删除、切换等基础管理功能，支持多快照并行维护，适配多场景实验需求。

\subsection{仿真场景二维快速预览模块}

仿真场景二维快速预览模块用于解决人车模拟器脚本开发中场景配置抽象、三维渲染依赖高、调试效率低的问题，实现脚本场景的轻量化可视化预览，实时响应代码变更，快速校验场景布局合理性。模块无需依赖三维渲染引擎，通过二维俯视图直观展示车辆、传感器、障碍物及行人的空间分布，降低场景调试门槛。

模块内部划分为代码解析、数据结构化、二维渲染、异常告警、实时联动五个子单元。代码解析子单元实时监听文档变更，通过抽象语法树解析技术提取脚本中create\_vehicle、add\_sensor、add\_pedestrian、add\_obstacle等核心指令的坐标、朝向及挂载关系；数据结构化子单元将解析结果整理为标准化对象列表，建立车辆与传感器的挂载关联；二维渲染子单元基于
WebView 与 Canvas
技术绘制俯视图，采用差异化图标区分车辆、传感器、障碍物等元素；异常告警子单元识别未定义车辆引用、坐标越界、传感器挂载异常等问题并高亮提示；实时联动子单元在代码修改后即时刷新预览界面，实现编码与预览的同步联动。

\subsection{仿真实验报告自动生成模块}

仿真实验报告自动生成模块用于解决人车模拟器实验记录零散、报告编写繁琐、格式不统一、数据整合效率低的问题，实现实验信息的一键采集、标准化整理与多格式导出，支撑实验成果的规范化归档与快速复盘。模块整合脚本信息、场景配置、设备参数、运行指标等全维度数据，生成符合科研与工程规范的实验报告。

模块内部划分为信息采集、数据整合、模板渲染、格式导出、报告管理五个子单元。信息采集子单元遍历仿真脚本与运行日志，自动提取脚本路径、场景文件、车辆与传感器清单、实验参数、批量测试用例及HUTB\_RESULT输出指标；数据整合子单元对采集的多源数据进行清洗、分类与结构化处理，形成标准化实验数据集；模板渲染子单元加载内置
Markdown 与 Word
格式报告模板，自动填充实验信息、数据图表与结论内容；格式导出子单元根据用户选择生成对应格式报告，统一存储于hutb-reports目录；报告管理子单元支持报告命名、预览、删除与批量导出，提升实验成果管理效率。

\subsection{常规仿真场景模块}

预制模板部分，本模块筛选出实验室最常用的三类场景：基础单车辆控制、多车协同和传感器数据采集，每个模板都包含完整的hutb.init\_simulator()、hutb.load\_scene()等标准化流程代码和详细注释，用户只需修改少量参数就能直接运行。

自定义模板部分，该模块设计了全局与工作区分开的存储机制，系统模板存在插件全局目录，用户自定义模板存在当前工作区目录，这样不同项目的模板不会混淆。模板文件采用
JSON 格式，结构简单易维护，无需额外解析库。

\subsection{语法高亮增强模块}

语法高亮增强模块用于解决通用编辑器无法识别人车模拟器专属
API、关键字与语法结构，导致代码可读性差、领域代码与通用代码区分模糊的问题，通过定制化语法规则实现
HUTB/MCP 专属代码元素的差异化着色，提升代码辨识度与开发效率。

模块内部划分为语法规则定义、语法注册、着色渲染、元素识别、高亮联动五个子单元。语法规则定义子单元基于
TextMate
语法规范，编写hutb.tmLanguage.json配置文件，定义hutb.、mcp.前缀、核心
API、导入语句、常量、关键字等元素的匹配规则；语法注册子单元在插件初始化阶段，将定制化语法注册至
VS Code 语言服务，关联 Python
文件类型；着色渲染子单元根据语法规则，为不同类型代码元素配置差异化颜色，区分关键字、API、模块名、常量、注释等；元素识别子单元实时解析代码文本，精准匹配专属语法元素；高亮联动子单元在代码编辑过程中实时更新高亮效果，保证着色与代码内容同步一致。

\section{系统接口设计}

代码补全接口：代码补全接口是编辑器和补全模块之间的接口，主要是监听用户输入，并给出相应的匹配结果。接口接收当前文档及光标的当前位置信息，利用前缀匹配的方式来快速定位到对应的功能API，并返回一个包含函数名、参数结构、功能描述、使用示例和代码片段的结构化数据包，从而保证了系统的响应速度为毫秒级，使得用户的输入过程保持流畅，同时也可以得到及时且准确的提示。

错误检测接口：错误检测接口接收到文档文本之后，利用静态语法分析、参数校验以及相似度匹配等手段，对代码里的语法错误、接口调用错误、参数不匹配、废弃函数等进行即时扫描，给出错误所在位置、类型、严重程度以及修正意见，从而达成编码阶段的即时提示，提前阻止运行时出现异常。

调试配置接口：调试配置接口自动生成调试配置文件，收到工作区路径、模板类型后自动写入launch.json配置，注入环境变量、绑定到Python解释器，可以实现各种场景模板的切换，简化调试准备工作、提高调试效率。

热重载接口：热重载接口监听文件保存事件，解析参数修改内容并进行合法性校验，用MCP协议向仿真引擎下发更新指令，返回执行状态和结果，实现仿真参数动态更新，不需要重启就可以生效。

监控面板接口：监控面板接口定时获取遥测数据，解析车速、加速度、位置、传感器状态等信息后返回结构化的数据，给前端可视化渲染提供支持，可以实现实时刷新、数据导出、离线模拟展示。

批量测试接口：批量测试接口读取测试配置文件，自动安排仿真脚本执行，收集运行日志，统计完成率，碰撞次数，耗时等指标，得到可视化测试报告，可以实现自动化验证和结果分析。

状态栏接口：状态栏接口可以实时得到系统的运行状态，显示Python版本、SDK以及系统是否连接、热重载的结果等信息，并且可以实现对于调试、打包等命令的快捷调用，从而使得调试变得更加方便。

快照与续跑接口：接收脚本上下文数据，完成快照捕获、序列化存储；接收快照标识，实现快照加载、调试配置更新与断点续跑指令下发，返回快照操作状态与执行结果。

场景预览接口：接收文档变更事件与代码文本，完成语法解析与场景数据提取；返回结构化场景数据，支撑二维视图渲染，同步返回异常告警信息。

实验报告接口：接收脚本路径、运行日志与测试数据，完成实验信息采集与整合；接收格式选择指令，生成对应格式报告并返回存储路径，支持报告预览与导出。

语法高亮接口：接收代码文本与语法规则配置，完成专属语法元素识别；返回差异化着色数据，同步更新编辑器高亮效果，支持语法规则动态更新。

模板管理接口：接收模板生成指令，返回对应场景模板代码；接收自定义模板保存指令，完成模板存储与列表更新；接收模板调用指令，返回模板代码，支撑脚本快速生成。

环境打包接口：环境打包接口接收到打包参数之后，就会把依赖资源复制到打包结果中，并自动生成配置文件以及启动脚本，最后压缩归档成一个便于分发和部署的安装包。

\section{系统运行流程}

系统采取事件驱动、模块化协同运转模式，全部运行流程包含插件启动、编码交流、校验调试、监控检测、打包部署全流程，各个模块按照同样的事件调度规则依次执行，具有反应迅速、耦合程度低、可扩展等优点，总体运行流程见图\ref{fig:3-3}。

插件初始化就是VS
Code启动之后会自动调用插件，读取API配置文件，完成各个功能模块的初始化、命令和事件的注册、环境依赖的检测以及状态栏的初始化，为之后的交互提供稳定的支撑。

编码交互中用户输入
hutb.或者mcp.前缀时，会触发代码补全模块对接口进行匹配并给出提示，错误检测模块也会对代码进行实时扫描，发现异常调用、参数错误等并加以标记，可以悬浮查看接口文档和示例。

调试配置中用户选择对应的调试模板，系统会自动生成launch.json配置文件，将环境变量注入进去，并与Python解释器进行匹配，从而达到一键启动调试的效果。

参数热重载就是用户对仿真参数进行修改并保存文件，系统会自动解析参数、合法性校验后用MCP协议发送到仿真引擎，状态栏实时显示执行结果。

实时监控，用户打开监控面板，系统定时获取仿真遥测数据，进行可视化图表渲染，可以实现数据实时展示、历史记录导出和异常告警。

批量测试是将用户的加载测试配置文件提交到系统中，系统就会自动调度仿真脚本来执行批量的测试任务，并且可以分析出执行的结果以及相关的指标数据，最后生成可视化的HTML测试报告。

环境打包时用户发出打包命令，系统就会把依赖资源复制下来、创建出配置和启动脚本，然后压缩成ZIP便携包。

整套流程形成一个完整的闭环，模块间相互配合、事件响应及时、操作链路简短，使得人车模拟器的开发过程一直处于稳定的状态之中。

\begin{figure}[H]
\centering
\begin{center}
\includegraphics[width=0.86\textwidth,height=0.58\textheight,keepaspectratio]{images/image6.png}
\end{center}
\caption{系统运行流程图}
\label{fig:3-3}
\end{figure}

\section{本章小结}

本章对系统的总体设计原则、分层架构、核心模块及接口设计、运行流程进行了详细阐述，新增场景快照与断点续跑、仿真场景二维快速预览、仿真实验报告自动生成、领域语义诊断、批量仿真测试增强、语法高亮增强、一键模板生成与自定义七大功能模块，进一步完善了系统功能矩阵，强化了系统在仿真回溯、场景可视化、实验规范化、错误精准检测、测试高效执行、代码可读性、脚本快速开发等方面的能力。模块化、低耦合的设计架构，保证了系统功能的完整性、可扩展性与稳定性，为后续系统实现与迭代优化提供了清晰的架构支撑。

